\caption{External sequence preferences and source-resolved natural variants. \textbf{a}, Two metric modules show all 25 proteins' absolute Frequency/Full scores above their paired Full-minus-Frequency differences: primary all-word Spearman and secondary composition-matched concordance. Narrow violins summarize the actual 25 scores per arm; points are individual proteins and short horizontal ticks are the original arm means. In the difference strips, grey points are the actual 25 paired differences, and red points and interval bars give the frozen paired-protein means and 95\% bootstrap intervals. The mean differences are $+0.146$ (95\% CI $+0.038$ to $+0.248$) for Spearman and $+3.78$ percentage points ($+0.90$ to $+6.73$) for composition-matched concordance. Complete Sequence controls remain in Figure~S31. \textbf{b}, All 30 QKI A$_2$C$_2$U words and 48 adjacent exchanges. Height encodes RNAcompete $Z_5$; node halves encode Full/Sequence mean responses across 32 computational backgrounds on the same shared asinh colour scale (navy--white--red). Three deep-blue edges identify CUAAC neighbours. \textbf{c}, Exact \textit{ACVR1B} A$>$G, source UPF1 preference for A (ASB FDR 0.0446), seven saved model/control responses and published HeLa RNA reporter measurements from \citet{yang2019asb}. Hollow shapes identify three fixed-input permutations, not uncertainty or biological replicates. Solid deep-blue reporter points are three biological measurements per allele with mean ticks; RNA abundance and predicted K562 binding are distinct endpoints. This is a post hoc training-locus illustration. \textbf{d}, All 396 development-exposed ASB-positive events (85 RBPs, 287 sites), with source-preferred-minus-other responses and five-arm RBP-equal accuracy/95\% intervals. Red/cyan identify Full-only/Sequence-only direction agreement; grey diagonal is $y=x$. \textbf{e}, All eight ClinVar P/LP and/or TCGA-overlap sites, 13 events, 52 effects and four-arm REF-minus-ALT responses, aligned by site and event. Light alternating backgrounds identify actual site groups. Positive/negative model-coloured bars extend right/left of the common zero line and favour REF/ALT; every source label favours REF. Hollow left/right triangles denote only the sign of nonzero bars shorter than 0.3 mm at the plotted size; raw bar lengths are unchanged. Agreement counts are 7/10/12/9 for Full/Sequence/flat/permuted input; the constant-REF rule gives 13/13. H/K denote HepG2/K562. These events include 11 training and two validation observations, with no strict-test events. Disease annotations do not establish binding truth or clinical causality. Both axes in d and all four effect axes in e share asinh width 0.02 within their panel; ticks retain original log-count units. All four e axes use the same range $[-0.12,0.88]$. The reference conservation input is fixed across each allele pair. Exact site identities: Figure~S32. Methods S1.18/S1.20--S1.22; Data S23/S25--S27.}
